\subsection{基本対称式型}\label{節：基本対称式型}
導入に続いて,先程の\bookterm{recurrence}{漸化式}を考えていこう.
\[
  a_1=3, \quad a_{n+1}=a_n^2-2
\]
この\bookterm{recurrence}{漸化式}を解くに当たって邪魔をしているのは$-2$の項である.
これをうまく消せるような操作を考えたい.
また,\,2乗があることからうまく因数分解できないか考えてみよう.

この$-2$を打ち消す$2$が生まれるような対称式を考える.

2乗の項は正だから,
\[
  (a+b)^2=a^2+2ab+b^2
\]
で言うところの$2ab$が$2$になってくれるように調整すればうまくいきそうだ.
そのようなふうに考えていくと,
\[
  \left(\alpha+\frac{1}{\alpha}\right)^2=\alpha^2+2+\frac{1}{\alpha^2}
\]
という式が思い当たる.
実際,ここに$-2$をすると
\[
  \alpha^2+\frac{1}{\alpha^2}
\]
が残るから,成り立ちそうなことが分かる.
この実験を元に解答に移る.

\noindent \bookblocklink{example-symmetric-1}{problem}{answer}{【問題】}\par\nobreak
\begin{smallquote}
  次の\bookterm{recurrence}{漸化式}で定義される\bookterm{sequence}{数列}の\bookterm{general}{一般項}を求めよ.
  \[
    a_1=3, \quad a_{n+1}=a_n^2-2
  \]
\end{smallquote}
\par\noindent\bookblocklink{example-symmetric-1}{answer}{problem}{【解答】}\par\nobreak
  \begin{smallquote}
  \bookterm{general}{一般項}が
    \[
      b_n=\alpha^{2^{n-1}}+\alpha^{-2^{n-1}}
    \]
    で表される\bookterm{sequence}{数列}$b_n$を考える.
    \begin{align*}
      b_{n+1}&=\alpha^{2^{n}}+\alpha^{-2^{n}}\\
      &=\left(\alpha^{2^{n-1}}\right)^2+\left(\alpha^{-2^{n-1}}\right)^2\\
      &=\left(\alpha^{2^{n-1}}+\alpha^{-2^{n-1}}\right)^2-2\\
      &=b_n^2-2
    \end{align*}
    これは与式の\bookterm{recurrence}{漸化式}と一致する.また,
    \[
      b_1=a_1=\alpha+\alpha^{-1}
    \]
    となる$\alpha$を考えると,
    \begin{align*}
      &\alpha+\alpha^{-1}=3\\
      &\alpha^2-3\alpha+1=0\\
      &\therefore \quad \alpha=\frac{3\pm \sqrt{5}}{2}
    \end{align*}
    得られた$\alpha$の2解は逆数の関係にある.
    以上より,
    \[
      a_n=\left(\frac{3+ \sqrt{5}}{2}\right)^{2^{n-1}}+\left(\frac{3- \sqrt{5}}{2}\right)^{2^{n-1}}
    \]
  \end{smallquote}

ここで,定数を $-2$ から $-1$ に変えてみよう.
\[
  a_{n+1}=a_n^2-1
\]
同じ置換を試すと,
\begin{align*}
  a_n^2-1
    &=\left(u+u^{-1}\right)^2-1\\
    &=u^2+u^{-2}+1,
\end{align*}
となる.しかし,次の項として欲しい形は $u^2+u^{-2}$ である.
余分な $1$ が残るため,同じ対称式の族の中に戻ってこない.

\bookterm{recurrence}{漸化式}の $-2$ と $-1$ は,見た目には小さな違いである.
しかし,対称式を保つ恒等式にとっては,その違いが決定的になる.
「ほとんど同じ式だから,ほとんど同じ方法で解ける」とは限らない.
解法が成立するかどうかは,係数の大きさではなく,恒等式が完全に一致
するかどうかで決まる.

同じ考え方は,三次式にも現れる.
\[
  \left(\alpha+\frac{1}{\alpha}\right)^3
  -3\left(\alpha+\frac{1}{\alpha}\right)
  =\alpha^3+\frac{1}{\alpha^3}
\]
このように,対称な組み合わせを選ぶと,べき乗の複雑な展開が
再び同じ種類の式へ戻る.ただし,恒等式が少しでもずれると,その閉じた
構造は失われる.これが基本対称式型の面白さである.

三乗の恒等式に合う\bookterm{recurrence}{漸化式}も作れる。

\noindent\bookblocklink{example-symmetric-2}{problem}{answer}{【問題】}\par\nobreak
\begin{smallquote}
  $\displaystyle a_1=\frac{8}{3},\quad a_{n+1}=a_n^3+3a_n$
\end{smallquote}
\par\noindent\bookblocklink{example-symmetric-2}{answer}{problem}{【解答】}\par\nobreak
\begin{smallquote}
  \bookterm{general}{一般項}が
    \[
      b_n=\alpha^{3^{n-1}}-\alpha^{-3^{n-1}}
    \]
    で表される\bookterm{sequence}{数列}$b_n$を考える.
    \begin{align*}
      b_{n+1}&=\alpha^{3^{n}}-\alpha^{-3^{n}}\\
      &=\left(\alpha^{3^{n-1}}\right)^3-\left(\alpha^{-3^{n-1}}\right)^3\\
      &=\left(\alpha^{3^{n-1}}-\alpha^{-3^{n-1}}\right)^3\\
      &\qquad +3\left(\alpha^{3^{n-1}}-\alpha^{-3^{n-1}}\right)\\
      &=b_n^3+3b_n
    \end{align*}
    これは与式の\bookterm{recurrence}{漸化式}と一致する.また,
    \[
      b_1=a_1=\alpha-\alpha^{-1}=\frac{8}{3}
    \]
    となる$\alpha$を考えると,
    \begin{align*}
      &\alpha-\alpha^{-1}=\frac{8}{3}\\
      &3\alpha^2-8\alpha-3=0\\
      &\therefore \quad \alpha=3,\,-\frac13
    \end{align*}
    どちらを選んでも結果は変わらないので,\,$\alpha=3$とすると,
    \[
      a_n=3^{3^{n-1}}-3^{-3^{n-1}}
    \]
\end{smallquote}
